% topics/second_order_map/main.tex
% Copyright (C) 2026 Hirokazu Maesaka (RIKEN SPring-8 Center)
% SPDX-License-Identifier: CC-BY-4.0
%
% Bilingual memo. Build with the commands in AGENTS.md (English: \memolang=en, Japanese: \memolang=ja).

% common/tex/preamble.tex
% Copyright (C) 2026 Hirokazu Maesaka (RIKEN SPring-8 Center)
% SPDX-License-Identifier: CC-BY-4.0
%
% Shared preamble for all topic documents. Loaded by topics/<topic>/main.tex with \input.
% Build with lualatex; the language is selected with \memolang (en or ja) through latexmk -pretex.

\documentclass[11pt,a4paper]{article}

% Language switch. The default is English. latexmk -pretex='\def\memolang{ja}' overrides it.
\providecommand{\memolang}{en}
% \ifx compares macro meanings. Unlike \ifdefstring, it is expandable, so \en and \ja also work inside \edef.
\def\memolangEN{en}
\def\memolangJA{ja}
\newcommand{\en}[1]{\ifx\memolang\memolangEN #1\fi}
\newcommand{\ja}[1]{\ifx\memolang\memolangJA #1\fi}

\usepackage{luatexja-preset}
\usepackage{etoolbox}
\usepackage{amsmath,amssymb}
\usepackage[style=numeric,backend=biber,sorting=none]{biblatex}
\addbibresource{common/references/references.bib}
% Memo numbers (YYYY-NNN) of the topics, generated from common/documents.toml. Cite a memo with \cite{YYYY-NNN}.
\addbibresource{common/documents.bib}

% Memo number. Each topic declares its number with \setmemonumber{YYYY-NNN} (see common/documents.toml).
% The number is shown in the page header and is used in the names of the output files.
\providecommand{\memonumber}{}
\newcommand{\setmemonumber}[1]{\def\memonumber{#1}}
\usepackage{fancyhdr}
\setlength{\headheight}{14pt}
\pagestyle{fancy}
\fancyhf{}
\fancyhead[R]{Memo~\memonumber}
\fancyfoot[C]{\thepage}
\fancypagestyle{plain}{%
  \fancyhf{}%
  \fancyhead[R]{Memo~\memonumber}%
  \fancyfoot[C]{\thepage}%
  \renewcommand{\headrulewidth}{0pt}}

\setlength{\parindent}{0pt}
\setlength{\parskip}{0.6em}

% Slide markers (write build/<topic>/<jobname>.slides for common/slides/build_slides.py).
% common/tex/slides.tex
% Copyright (C) 2026 Hirokazu Maesaka (RIKEN SPring-8 Center)
% SPDX-License-Identifier: CC-BY-4.0
%
% Slide markers. Each marker writes one line "kind:payload" to build/<topic>/<jobname>.slides during compilation.
% common/slides/build_slides.py reads that file and builds the PowerPoint deck. The markers print nothing in the PDF.
% Arguments are expanded at the point of use, so \en{...} and \ja{...} select the language of the current build.
% Do not use fragile commands (for example \cite or \label) inside slide arguments.

\newwrite\slidesfile
\AtBeginDocument{\immediate\openout\slidesfile=\jobname.slides\relax}
\AtEndDocument{\immediate\closeout\slidesfile}

% #1: entry kind, #2: payload. The payload is expanded, then written as plain characters.
\newcommand{\slideentry}[2]{%
  \begingroup
  \edef\slidetmp{#2}%
  \edef\slidedet{\detokenize\expandafter{\slidetmp}}%
  \immediate\write\slidesfile{#1:\slidedet}%
  \endgroup}

% Raw entries are written without expansion. Use them for math and paths, which contain fragile commands.
\newcommand{\slideentryraw}[2]{%
  \immediate\write\slidesfile{#1:\detokenize{#2}}}

\newcommand{\slidetitle}[1]{\slideentry{title}{#1}}
\newcommand{\slidepoint}[1]{\slideentry{point}{#1}}
\newcommand{\slidemath}[1]{\slideentryraw{math}{#1}}
\newcommand{\slidefigure}[1]{\slideentryraw{figure}{#1}}
\newcommand{\slidenote}[1]{\slideentry{note}{#1}}


\setmemonumber{2026-004}

\title{\en{Second-order transfer maps of a drift space, a sector dipole magnet, and a quadrupole magnet}%
       \ja{ドリフトスペース，扇型偏向磁石，四極磁石の 2 次の転送写像}}
\author{Hirokazu Maesaka}
\date{Memo~\memonumber, 2026}

\begin{document}
\maketitle

% Title slide. The slides of each section are written at the end of that section (see AGENTS.md).
\slidetitle{\en{Second-order transfer maps of a drift space, a sector dipole magnet, and a quadrupole magnet}%
            \ja{ドリフトスペース，扇型偏向磁石，四極磁石の 2 次の転送写像}}
\slidepoint{\en{Second-order terms T_ijk from the equations of motion through second order}%
            \ja{2 次までの運動方程式から 2 次の項 T_ijk を導く}}
\slidenote{\en{Starts from Memo 2026-002 (linear matrices) and Memo 2026-003 (second-order EOM).}%
           \ja{メモ 2026-002（線形行列）とメモ 2026-003（2 次の運動方程式）から出発する。}}

\section{\en{Scope and conventions}\ja{対象と記号の規約}}

\en{This memo derives the second-order transfer map of a drift space, a sector dipole magnet, and a quadrupole
magnet. The starting points are the linear transfer matrices of \cite{2026-002} and the equations of motion (EOM)
through second order of \cite[Sec.~3]{2026-003}. The assumptions are those of \cite{2026-002}: static field,
$k_s = 0$, constant $h$ and $k$ in each element, hard-edge model, and a sector dipole magnet with perpendicular end
faces. Quadrupole curvature corrections, which are neglected in \cite{2026-001}, are not included. Every result is
verified with sympy in the companion notebook \texttt{second\_order\_map.py}, and numerically against the exact
EOM.}%
\ja{本メモでは，ドリフトスペース，扇型偏向磁石，四極磁石の 2 次の転送写像を導く。出発点は，\cite{2026-002} の
線形転送行列と，\cite[第 3 節]{2026-003} の 2 次までの運動方程式である。仮定は \cite{2026-002} と同じで，
静磁場，$k_s = 0$，各要素の中で $h$ と $k$ が一定，ハードエッジモデル，端面が垂直な扇型偏向磁石である。
\cite{2026-001} で無視した四極磁石の曲率補正は含めない。すべての結果は，付随するノートブック
\texttt{second\_order\_map.py} で sympy により検証し，さらに厳密な運動方程式に対して数値的にも確認している。}

\en{The phase-space vector is $\mathbf{u} = (x, p_x, y, p_y, z, \delta)^\mathsf{T}$ with the conventions of
\cite{2026-001}. The momenta $p_x$ and $p_y$ are the canonical momenta normalized by $P_0$; they are not the slopes
$x'$ and $y'$, which differ from them at second order.}%
\ja{位相空間ベクトルは $\mathbf{u} = (x, p_x, y, p_y, z, \delta)^\mathsf{T}$ であり，記号の規約は \cite{2026-001}
に従う。運動量 $p_x$，$p_y$ は $P_0$ で規格化した正準運動量であり，傾き $x'$，$y'$ ではない。両者は 2 次で
異なる。}

% Slides for this section.
\slidetitle{\en{Scope and conventions}\ja{対象と記号の規約}}
\slidepoint{\en{Elements: drift space, sector dipole magnet, quadrupole magnet}%
            \ja{対象: ドリフトスペース，扇型偏向磁石，四極磁石}}
\slidepoint{\en{Starting points: Memo 2026-002 (R_ij) and Memo 2026-003 (EOM)}%
            \ja{出発点: メモ 2026-002（R_ij）とメモ 2026-003（運動方程式）}}
\slidepoint{\en{p_x, p_y are canonical momenta normalized by P0, not slopes}%
            \ja{p_x，p_y は P0 で規格化した正準運動量であり，傾きではない}}

\section{\en{Second-order solution of the equations of motion}\ja{運動方程式の 2 次の解}}
\label{sec:duhamel}

\en{The EOM through second order of \cite[Sec.~3]{2026-003} are $\mathbf{u}' = \mathbf{A}\mathbf{u} +
\mathbf{f}_2(\mathbf{u})$, where $\mathbf{A}$ is the matrix of \cite[Sec.~2]{2026-002} and $\mathbf{f}_2$ is the
quadratic part obtained from the cubic Hamiltonian $H_3$,}%
\ja{2 次までの運動方程式 \cite[第 3 節]{2026-003} は $\mathbf{u}' = \mathbf{A}\mathbf{u} +
\mathbf{f}_2(\mathbf{u})$ である。ここで $\mathbf{A}$ は \cite[第 2 節]{2026-002} の行列であり，$\mathbf{f}_2$ は
3 次のハミルトニアン $H_3$ から得られる 2 次の部分である。}
\begin{equation}
  H_3 = (hx - \delta)\,\frac{p_x^2 + p_y^2}{2} + (hx - \delta)\,\frac{\delta^2}{2\gamma_0^2},
  \label{eq:H3}
\end{equation}
\begin{equation}
  \mathbf{f}_2 = \Big(\frac{\partial H_3}{\partial p_x},\ -\frac{\partial H_3}{\partial x},\
  \frac{\partial H_3}{\partial p_y},\ -\frac{\partial H_3}{\partial y},\ \frac{\partial H_3}{\partial \delta},\
  0\Big)^\mathsf{T} .
  \label{eq:f2}
\end{equation}
\en{Write the solution as $\mathbf{u}(L) = \mathbf{M}(L)\,\mathbf{u}_0 + \mathbf{w}(L)$, where $\mathbf{M}$ is the
linear transfer matrix and $\mathbf{w}$ is of second order in the entrance vector $\mathbf{u}_0$. Inserting this
into the EOM and keeping the second-order terms, where $\mathbf{f}_2$ is evaluated on the first-order solution
$\mathbf{M}\mathbf{u}_0$, gives}%
\ja{解を $\mathbf{u}(L) = \mathbf{M}(L)\,\mathbf{u}_0 + \mathbf{w}(L)$ と書く。$\mathbf{M}$ は線形転送行列，
$\mathbf{w}$ は入口のベクトル $\mathbf{u}_0$ について 2 次の項である。これを運動方程式に代入し，
$\mathbf{f}_2$ を 1 次の解 $\mathbf{M}\mathbf{u}_0$ で評価して 2 次の項を残すと}
\begin{equation}
  \mathbf{w}' = \mathbf{A}\,\mathbf{w} + \mathbf{f}_2\big(\mathbf{M}(L)\,\mathbf{u}_0\big), \qquad \mathbf{w}(0) = 0 .
  \label{eq:w_ode}
\end{equation}
\en{With $\mathbf{w} = \mathbf{M}\mathbf{v}$ and $\mathbf{M}' = \mathbf{A}\mathbf{M}$ this becomes
$\mathbf{v}' = \mathbf{M}^{-1}\mathbf{f}_2$, and $\mathbf{M}(L)\,\mathbf{M}(s)^{-1} = \mathbf{M}(L - s)$ for a
constant $\mathbf{A}$. The variation-of-constants (Duhamel) solution is therefore}%
\ja{$\mathbf{w} = \mathbf{M}\mathbf{v}$ とおき $\mathbf{M}' = \mathbf{A}\mathbf{M}$ を用いると
$\mathbf{v}' = \mathbf{M}^{-1}\mathbf{f}_2$ となる。$\mathbf{A}$ が一定なら
$\mathbf{M}(L)\,\mathbf{M}(s)^{-1} = \mathbf{M}(L - s)$ である。したがって定数変化法（Duhamel の公式）による解は}
\begin{equation}
  \mathbf{w}(L) = \int_0^L \mathbf{M}(L - s)\;\mathbf{f}_2\big(\mathbf{M}(s)\,\mathbf{u}_0\big)\,ds .
  \label{eq:duhamel}
\end{equation}
\en{Each component of $\mathbf{w}$ is a quadratic form in $\mathbf{u}_0$. The second-order transfer map is defined by}%
\ja{$\mathbf{w}$ の各成分は $\mathbf{u}_0$ の 2 次形式である。2 次の転送写像は}
\begin{equation}
  u_i(L) = \sum_j R_{ij}\,u_{0j} + \sum_{j \le k} T_{ijk}\,u_{0j}\,u_{0k}, \qquad R_{ij} = M_{ij},
  \label{eq:Tijk}
\end{equation}
\en{so $T_{ijk}$ is the coefficient of the monomial $u_{0j}u_{0k}$ in $w_i$, with each pair $j \le k$ counted once
(the convention of TRANSPORT). The notebook evaluates \eqref{eq:duhamel} with sympy for each element, one monomial
at a time. All integrands are products of $\sin$, $\cos$, $\sinh$, $\cosh$, and powers of $s$.}%
\ja{により定義する。$T_{ijk}$ は $w_i$ における単項式 $u_{0j}u_{0k}$ の係数であり，各組 $j \le k$ を 1 回だけ数える
（TRANSPORT の規約）。ノートブックでは，各要素について \eqref{eq:duhamel} を単項式ごとに sympy で積分する。
被積分関数は $\sin$，$\cos$，$\sinh$，$\cosh$ と $s$ の冪の積である。}

% Slides for this section.
\slidetitle{\en{Second-order solution}\ja{運動方程式の 2 次の解}}
\slidemath{\mathbf{u}' = \mathbf{A}\mathbf{u} + \mathbf{f}_2(\mathbf{u}), \qquad \mathbf{u}(L) = \mathbf{M}\mathbf{u}_0 + \mathbf{w}(L)}
\slidemath{\mathbf{w}' = \mathbf{A}\,\mathbf{w} + \mathbf{f}_2\big(\mathbf{M}(L)\,\mathbf{u}_0\big), \qquad \mathbf{w}(0) = 0}
\slidemath{\mathbf{w}(L) = \int_0^L \mathbf{M}(L - s)\;\mathbf{f}_2\big(\mathbf{M}(s)\,\mathbf{u}_0\big)\,ds}

\slidetitle{\en{Second-order transfer map}\ja{2 次の転送写像}}
\slidemath{u_i(L) = \sum_j R_{ij}\,u_{0j} + \sum_{j \le k} T_{ijk}\,u_{0j}\,u_{0k}}
\slidepoint{\en{T_ijk: coefficient of u_0j u_0k in w_i, each pair j <= k once}%
            \ja{T_ijk は w_i における u_0j u_0k の係数（j ≤ k を 1 回ずつ）}}

\section{\en{Drift space}\ja{ドリフトスペース}}
\label{sec:drift2}

\en{With $h = 0$ the only second-order terms of $\mathbf{f}_2$ are $x' \ni -\delta\,p_x$, $y' \ni -\delta\,p_y$, and
$z' \ni -(p_x^2 + p_y^2)/2 - 3\delta^2/(2\gamma_0^2)$, and the first-order solution has constant $p_x$, $p_y$, and
$\delta$. Integrating over $L$ gives}%
\ja{$h = 0$ では $\mathbf{f}_2$ の 2 次の項は $x' \ni -\delta\,p_x$，$y' \ni -\delta\,p_y$，
$z' \ni -(p_x^2 + p_y^2)/2 - 3\delta^2/(2\gamma_0^2)$ だけであり，1 次の解では $p_x$，$p_y$，$\delta$ は一定である。
$L$ にわたって積分すると}
\begin{equation}
  \begin{aligned}
    w_1 &= -L\,\delta_0\,p_{x0}, \qquad w_3 = -L\,\delta_0\,p_{y0}, \qquad w_2 = w_4 = w_6 = 0, \\
    w_5 &= -\frac{L}{2}\,(p_{x0}^2 + p_{y0}^2) - \frac{3L}{2\gamma_0^2}\,\delta_0^2 .
  \end{aligned}
  \label{eq:w_drift}
\end{equation}
\en{The nonzero elements are $T_{126}$, $T_{346}$, $T_{522}$, $T_{544}$, and $T_{566}$. The terms $-L\delta_0 p_{x0}$
and $-L\delta_0 p_{y0}$ are the change of the slope with the momentum, $x' = p_x/(1 + \delta) \approx p_x(1 - \delta)$,
and $-L(p_{x0}^2 + p_{y0}^2)/2$ is the path-length increase of a slanted trajectory.}%
\ja{非零の要素は $T_{126}$，$T_{346}$，$T_{522}$，$T_{544}$，$T_{566}$ である。$-L\delta_0 p_{x0}$ と
$-L\delta_0 p_{y0}$ は，運動量による傾きの変化 $x' = p_x/(1 + \delta) \approx p_x(1 - \delta)$ であり，
$-L(p_{x0}^2 + p_{y0}^2)/2$ は，傾いた軌道の経路長の増加である。}

% Slides for this section.
\slidetitle{\en{Drift space: second-order terms}\ja{ドリフトスペース: 2 次の項}}
\slidemath{w_1 = -L\,\delta_0\,p_{x0}, \qquad w_3 = -L\,\delta_0\,p_{y0}}
\slidemath{w_5 = -\frac{L}{2}\,(p_{x0}^2 + p_{y0}^2) - \frac{3L}{2\gamma_0^2}\,\delta_0^2}
\slidepoint{\en{w_2 = w_4 = w_6 = 0; nonzero: T126, T346, T522, T544, T566}%
            \ja{w_2 = w_4 = w_6 = 0。非零: T126，T346，T522，T544，T566}}

\section{\en{Sector dipole magnet}\ja{扇型偏向磁石}}
\label{sec:bend2}

\en{With $h \neq 0$ and $k = 0$, let $\theta = hL$. The first-order solution of \cite[Sec.~5]{2026-002} is
inserted into $\mathbf{f}_2$, and \eqref{eq:duhamel} contains only products of $\sin hs$ and $\cos hs$. The result
is}%
\ja{$h \neq 0$，$k = 0$ とし，$\theta = hL$ とおく。\cite[第 5 節]{2026-002} の 1 次の解を $\mathbf{f}_2$ に代入すると，
\eqref{eq:duhamel} の被積分関数は $\sin hs$ と $\cos hs$ の積だけになる。結果は}
\begin{equation}
  \begin{aligned}
    w_1 &= -\frac{h\sin^2\theta}{2}\,x_0^2 + \sin\theta\cos\theta\,x_0 p_{x0} + \sin^2\theta\,x_0\delta_0
      - \frac{\sin\theta\cos\theta}{h}\,p_{x0}\delta_0 \\
      &\quad + \frac{\cos\theta\,(1-\cos\theta)}{2h}\,p_{x0}^2 - \frac{1-\cos\theta}{2h}\,p_{y0}^2
      - \Big[\frac{\sin^2\theta}{2h} + \frac{1-\cos\theta}{2\gamma_0^2 h}\Big]\delta_0^2 , \\
    w_2 &= -\frac{\sin\theta}{2}\Big(p_{x0}^2 + p_{y0}^2 + \frac{\delta_0^2}{\gamma_0^2}\Big), \\
    w_3 &= \sin\theta\,x_0 p_{y0} + \frac{1-\cos\theta}{h}\,p_{x0}p_{y0} - \frac{\sin\theta}{h}\,p_{y0}\delta_0 , \\
    w_5 &= -\frac{\sin\theta}{2h}\,(p_{x0}^2 + p_{y0}^2) + \frac{\sin\theta}{\gamma_0^2}\,x_0\delta_0
      + \frac{1-\cos\theta}{\gamma_0^2 h}\,p_{x0}\delta_0 - \frac{3\sin\theta}{2\gamma_0^2 h}\,\delta_0^2 , \\
    w_4 &= w_6 = 0 .
  \end{aligned}
  \label{eq:w_bend}
\end{equation}
\en{There are 18 nonzero elements $T_{ijk}$. The terms in $w_2$ are the sagitta effect: a particle with a transverse
slope or a momentum deviation travels on a longer arc in the bend, so $p_x$ changes by $-\sin\theta\,p_\perp^2/2$. The
terms of $w_1$ with $\delta_0$ are the chromatic changes of the orbit, and the terms of $w_5$ are the path-length
changes. For $h \to 0$ (with $L$ fixed) \eqref{eq:w_bend} tends to \eqref{eq:w_drift}.}%
\ja{非零の要素 $T_{ijk}$ は 18 個ある。$w_2$ の項は矢高の効果である。横方向に傾いた粒子や運動量偏差を持つ粒子は
偏向部でより長い弧を通るので，$p_x$ は $-\sin\theta\,p_\perp^2/2$ だけ変化する。$w_1$ の $\delta_0$ を含む項は
軌道の運動量による変化であり，$w_5$ の項は経路長の変化である。$h \to 0$（$L$ 固定）で \eqref{eq:w_bend} は
\eqref{eq:w_drift} に近づく。}

% Slides for this section.
\slidetitle{\en{Sector dipole magnet: x}\ja{扇型偏向磁石: x}}
\slidepoint{\en{theta = hL; Duhamel integrals contain only sin and cos of hs}%
            \ja{θ = hL。Duhamel の積分は sin hs と cos hs の積だけ}}
\slidemath{w_1 = -\frac{h\sin^2\theta}{2}\,x_0^2 + \sin\theta\cos\theta\,x_0 p_{x0} + \sin^2\theta\,x_0\delta_0}
\slidemath{\qquad - \frac{\sin\theta\cos\theta}{h}\,p_{x0}\delta_0 + \frac{\cos\theta\,(1-\cos\theta)}{2h}\,p_{x0}^2}
\slidemath{\qquad - \frac{1-\cos\theta}{2h}\,p_{y0}^2 - \Big[\frac{\sin^2\theta}{2h} + \frac{1-\cos\theta}{2\gamma_0^2 h}\Big]\delta_0^2}

\slidetitle{\en{Sector dipole magnet: p_x, y, and z}\ja{扇型偏向磁石: p_x，y，z}}
\slidemath{w_2 = -\frac{\sin\theta}{2}\Big(p_{x0}^2 + p_{y0}^2 + \frac{\delta_0^2}{\gamma_0^2}\Big)}
\slidemath{w_3 = \sin\theta\,x_0 p_{y0} + \frac{1-\cos\theta}{h}\,p_{x0}p_{y0} - \frac{\sin\theta}{h}\,p_{y0}\delta_0}
\slidemath{w_5 = -\frac{\sin\theta}{2h}\,(p_{x0}^2 + p_{y0}^2) + \frac{\sin\theta}{\gamma_0^2}\,x_0\delta_0}
\slidemath{\qquad + \frac{1-\cos\theta}{\gamma_0^2 h}\,p_{x0}\delta_0 - \frac{3\sin\theta}{2\gamma_0^2 h}\,\delta_0^2}

\section{\en{Quadrupole magnet}\ja{四極磁石}}
\label{sec:quad2}

\en{With $h = 0$ and $k = g^2 > 0$, let $\varphi = gL$. In $\mathbf{f}_2$ only the terms $x' \ni -\delta p_x$,
$y' \ni -\delta p_y$ and the $z'$ terms of the drift space remain, so the second-order terms are the effect of
these terms on the first-order oscillation \cite[Sec.~6]{2026-002}:}%
\ja{$h = 0$，$k = g^2 > 0$ とし，$\varphi = gL$ とおく。$\mathbf{f}_2$ に残るのは $x' \ni -\delta p_x$，
$y' \ni -\delta p_y$ とドリフトスペースの $z'$ の項だけであり，2 次の項はこれらが 1 次の振動
\cite[第 6 節]{2026-002} に及ぼす効果である。}
\begin{equation}
  \begin{aligned}
    w_1 &= \frac{\delta_0}{2g}\Big[g^2 L\sin\varphi\;x_0 - (\varphi\cos\varphi + \sin\varphi)\,p_{x0}\Big], \\
    w_2 &= \frac{g\,\delta_0}{2}\Big[(\varphi\cos\varphi - \sin\varphi)\,x_0 + L\sin\varphi\;p_{x0}\Big], \\
    w_3 &= -\frac{\delta_0}{2g}\Big[g^2 L\sinh\varphi\;y_0 + (\varphi\cosh\varphi + \sinh\varphi)\,p_{y0}\Big], \\
    w_4 &= -\frac{g\,\delta_0}{2}\Big[(\varphi\cosh\varphi - \sinh\varphi)\,y_0 + L\sinh\varphi\;p_{y0}\Big], \\
    w_5 &= -\frac{3L}{2\gamma_0^2}\,\delta_0^2
      - \frac{1}{8g}\Big[g^2(2\varphi - \sin 2\varphi)\,x_0^2 - 4g\sin^2\varphi\;x_0 p_{x0}
      + (2\varphi + \sin 2\varphi)\,p_{x0}^2 \\
      &\qquad - g^2(2\varphi - \sinh 2\varphi)\,y_0^2 + 4g\sinh^2\varphi\;y_0 p_{y0}
      + (2\varphi + \sinh 2\varphi)\,p_{y0}^2\Big], \\
    w_6 &= 0 .
  \end{aligned}
  \label{eq:w_quad}
\end{equation}
\en{There are 15 nonzero elements $T_{ijk}$. All second-order terms of $x$, $p_x$, $y$, and $p_y$ are proportional to
$\delta_0$, which is the chromatic effect of $x' = p_x(1 - \delta)$ on the oscillation. The terms of $w_5$ come from
the slopes and from $\delta_0^2$. For $k = -g^2 < 0$ the equations of the two planes are exchanged, so
$\mathbf{w}_{k<0}(\mathbf{u}_0) = \mathbf{P}\,\mathbf{w}_{k>0}(\mathbf{P}\,\mathbf{u}_0)$, where the permutation
$\mathbf{P}$ exchanges $(x, p_x)$ and $(y, p_y)$.}%
\ja{非零の要素 $T_{ijk}$ は 15 個ある。$x$，$p_x$，$y$，$p_y$ の 2 次の項はすべて $\delta_0$ に比例する。これは
$x' = p_x(1 - \delta)$ が振動に及ぼすクロマティックな効果である。$w_5$ の項は傾きと $\delta_0^2$ から生じる。
$k = -g^2 < 0$ では二つの面の方程式が入れ替わるので，$(x, p_x)$ と $(y, p_y)$ を入れ替える置換 $\mathbf{P}$ を用いて
$\mathbf{w}_{k<0}(\mathbf{u}_0) = \mathbf{P}\,\mathbf{w}_{k>0}(\mathbf{P}\,\mathbf{u}_0)$ である。}

\en{In the thin-lens limit, $L \to 0$ with $q = kL$ fixed, all of $\mathbf{w}$ tends to zero: a thin quadrupole has no
second-order terms in the canonical variables $x$ and $p_x$, because the kick $\Delta p_x = -q\,x$ does not depend on
$\delta$. The chromatic focusing appears only through the conversion $x' = p_x(1 - \delta)$ in the following drift
space.}%
\ja{薄肉レンズ近似（$q = kL$ を一定にして $L \to 0$）では $\mathbf{w}$ はすべてゼロに近づく。キック
$\Delta p_x = -q\,x$ は $\delta$ に依らないので，薄い四極磁石は正準変数 $x$，$p_x$ について 2 次の項を持たない。
クロマティックな集束は，続くドリフトスペースでの $x' = p_x(1 - \delta)$ への変換を通してだけ現れる。}

% Slides for this section.
\slidetitle{\en{Quadrupole magnet (k = g^2 > 0)}\ja{四極磁石（k = g^2 > 0）}}
\slidemath{w_1 = \frac{\delta_0}{2g}\Big[g^2 L\sin\varphi\;x_0 - (\varphi\cos\varphi + \sin\varphi)\,p_{x0}\Big]}
\slidemath{w_2 = \frac{g\,\delta_0}{2}\Big[(\varphi\cos\varphi - \sin\varphi)\,x_0 + L\sin\varphi\;p_{x0}\Big]}
\slidemath{w_3 = -\frac{\delta_0}{2g}\Big[g^2 L\sinh\varphi\;y_0 + (\varphi\cosh\varphi + \sinh\varphi)\,p_{y0}\Big]}
\slidemath{w_4 = -\frac{g\,\delta_0}{2}\Big[(\varphi\cosh\varphi - \sinh\varphi)\,y_0 + L\sinh\varphi\;p_{y0}\Big]}

\slidetitle{\en{Quadrupole magnet: z}\ja{四極磁石: z}}
\slidemath{w_5 = -\frac{3L}{2\gamma_0^2}\,\delta_0^2 - \frac{1}{8g}\Big[g^2(2\varphi - \sin 2\varphi)\,x_0^2}
\slidemath{\qquad - 4g\sin^2\varphi\;x_0 p_{x0} + (2\varphi + \sin 2\varphi)\,p_{x0}^2 + (y\ \text{terms})\Big]}
\slidepoint{\en{k < 0: exchange (x, p_x) and (y, p_y)}\ja{k < 0: (x, p_x) と (y, p_y) を入れ替える}}
\slidepoint{\en{Thin lens: all second-order terms vanish in the canonical variables}%
            \ja{薄肉レンズ近似: 正準変数では 2 次の項はすべて消える}}

\section{\en{Symplecticity of the truncated map}\ja{打ち切った写像のシンプレクティック性}}
\label{sec:symplectic}

\en{The exact map is symplectic. The map truncated at second order is not, and this section quantifies the defect.
Write $\mathbf{D} = \partial\mathbf{P}/\partial\mathbf{u}_0 = \mathbf{M} + \mathbf{W}$ with
$\mathbf{W} = \partial\mathbf{w}/\partial\mathbf{u}_0$. Since $\mathbf{M}$ is symplectic,
$\mathbf{M}^\mathsf{T}\mathbf{J}\mathbf{M} = \mathbf{J}$, the defect splits into two orders,}%
\ja{厳密な写像はシンプレクティックである。2 次で打ち切った写像はそうではなく，本節ではその欠陥を定量化する。
$\mathbf{D} = \partial\mathbf{P}/\partial\mathbf{u}_0 = \mathbf{M} + \mathbf{W}$，$\mathbf{W} = \partial\mathbf{w}/\partial\mathbf{u}_0$
とおく。$\mathbf{M}$ はシンプレクティックで $\mathbf{M}^\mathsf{T}\mathbf{J}\mathbf{M} = \mathbf{J}$ であるから，
欠陥は 2 つの次数に分かれる。}
\begin{equation}
  \mathbf{D}^\mathsf{T}\mathbf{J}\mathbf{D} - \mathbf{J}
  = \underbrace{\mathbf{M}^\mathsf{T}\mathbf{J}\mathbf{W} + \mathbf{W}^\mathsf{T}\mathbf{J}\mathbf{M}}_{\mathcal{O}(u_0)}
  + \underbrace{\mathbf{W}^\mathsf{T}\mathbf{J}\mathbf{W}}_{\mathcal{O}(u_0^2)} .
  \label{eq:defect_split}
\end{equation}
\en{The first group vanishes identically for every element: this is the second-order symplectic condition checked
below. The first non-vanishing term is therefore $\mathbf{W}^\mathsf{T}\mathbf{J}\mathbf{W}$, which is quadratic in
$\mathbf{u}_0$. The map truncated at second order is thus symplectic through first order in $\mathbf{u}_0$ and not
beyond.}%
\ja{第 1 群はすべての要素で恒等的にゼロである。これが以下で確認する 2 次のシンプレクティック条件である。
したがって最初に残る項は $\mathbf{W}^\mathsf{T}\mathbf{J}\mathbf{W}$ であり，$\mathbf{u}_0$ について 2 次である。
2 次で打ち切った写像は $\mathbf{u}_0$ の 1 次までシンプレクティックであり，それ以上ではない。}

\en{The result for the three elements is summarised in Table~\ref{tab:defect}.}%
\ja{3 つの要素についての結果を表~\ref{tab:defect} にまとめる。}
\begin{table}[htbp]
  \centering
  \begin{tabular}{lll}
    \hline
    \en{Element}\ja{要素} & $\mathbf{W}^\mathsf{T}\mathbf{J}\mathbf{W}$ & $\det\mathbf{D} - 1$ \\
    \hline
    \en{Drift space}\ja{ドリフトスペース} & $0$ (\en{exact}\ja{厳密}) & $0$ (\en{exact}\ja{厳密}) \\
    \en{Sector dipole}\ja{扇型偏向磁石} & $\sin^2\theta \times$ \en{rank 2}\ja{階数 2} & $-p_{x0}\,p_x(L)\sin^2\theta$ \\
    \en{Quadrupole}\ja{四極磁石} & $\delta_0^2 \times$ \en{rank 4}\ja{階数 4} & $\approx L^4\delta_0^2 g^4/6$ \\
    \hline
  \end{tabular}
  \caption{\en{Defect of the second-order map.}\ja{2 次の写像の欠陥。}}
  \label{tab:defect}
\end{table}

\en{\textbf{Drift space.} The truncated map is the exact flow of the polynomial Hamiltonian}%
\ja{\textbf{ドリフトスペース。}打ち切った写像は，多項式ハミルトニアン}
\begin{equation}
  K = (1-\delta)\,\frac{p_x^2 + p_y^2}{2} + \frac{\delta^2}{2\gamma_0^2} - \frac{\delta^3}{2\gamma_0^2}
  \label{eq:K_drift}
\end{equation}
\en{because $K$ does not depend on $x$, $y$, or $z$, so the momenta are constant and $x' = \partial K/\partial p_x$,
$y' = \partial K/\partial p_y$, $z' = \partial K/\partial\delta$ integrate directly. A flow of a Hamiltonian is
symplectic, so the drift map is symplectic to all orders, not only through second order.}%
\ja{の厳密な流れである。$K$ は $x$，$y$，$z$ に依らないので運動量は一定であり，
$x' = \partial K/\partial p_x$，$y' = \partial K/\partial p_y$，$z' = \partial K/\partial\delta$ はそのまま積分できる。
ハミルトニアンの流れはシンプレクティックであるから，ドリフトの写像は 2 次までではなくすべての次数で
シンプレクティックである。}

\en{\textbf{Sector dipole magnet.} The defect is proportional to $\sin^2\theta$, so it is small for a weak bend but
never vanishes. The determinant gives a compact measure,}%
\ja{\textbf{扇型偏向磁石。}欠陥は $\sin^2\theta$ に比例するので，弱い偏向では小さいがゼロにはならない。
行列式が簡潔な指標を与える。}
\begin{equation}
  \det\mathbf{D} - 1 = -p_{x0}\,p_x(L)\,\sin^2\theta ,
  \qquad p_x(L) = \delta_0\sin\theta - h x_0\sin\theta + p_{x0}\cos\theta .
  \label{eq:det_bend}
\end{equation}
\en{The defect vanishes on the reference orbit and for a particle with $p_{x0} = 0$ or $p_x(L) = 0$, and it is of
order $\theta$ in the $(y, p_y, \delta)$ block and of order $\theta^2$ in the $(x, p_x)$ block.}%
\ja{欠陥は基準軌道上で，また $p_{x0} = 0$ または $p_x(L) = 0$ の粒子でゼロになり，$(y, p_y, \delta)$ ブロックでは
$\theta$ の次数，$(x, p_x)$ ブロックでは $\theta^2$ の次数である。}

\en{\textbf{Quadrupole magnet.} With $A = (Lg)^2 - \sin^2(Lg)$ and $B = (Lg)^2 - \sinh^2(Lg)$, the twelve nonzero
entries of $\mathbf{W}^\mathsf{T}\mathbf{J}\mathbf{W}$ are}%
\ja{\textbf{四極磁石。}$A = (Lg)^2 - \sin^2(Lg)$，$B = (Lg)^2 - \sinh^2(Lg)$ とすると，
$\mathbf{W}^\mathsf{T}\mathbf{J}\mathbf{W}$ の非零成分 12 個は}
\begin{equation}
  \begin{aligned}
    \big(\mathbf{W}^\mathsf{T}\mathbf{J}\mathbf{W}\big)_{12} &= \frac{\delta_0^2 A}{4} , &
    \big(\mathbf{W}^\mathsf{T}\mathbf{J}\mathbf{W}\big)_{16} &= \frac{\delta_0 p_{x0} A}{4} , &
    \big(\mathbf{W}^\mathsf{T}\mathbf{J}\mathbf{W}\big)_{26} &= -\frac{\delta_0 x_0 A}{4} , \\
    \big(\mathbf{W}^\mathsf{T}\mathbf{J}\mathbf{W}\big)_{34} &= -\frac{\delta_0^2 B}{4} , &
    \big(\mathbf{W}^\mathsf{T}\mathbf{J}\mathbf{W}\big)_{36} &= -\frac{\delta_0 p_{y0} B}{4} , &
    \big(\mathbf{W}^\mathsf{T}\mathbf{J}\mathbf{W}\big)_{46} &= \frac{\delta_0 y_0 B}{4} ,
  \end{aligned}
  \label{eq:defect_quad}
\end{equation}
\en{with the antisymmetric partners. The defect is proportional to $\delta_0^2$, so it vanishes for a particle at the
reference momentum, and it vanishes in the thin-lens limit because $A, B = \mathcal{O}((Lg)^4)$ while the map itself
is of order $L$.}%
\ja{と反対称の成分である。欠陥は $\delta_0^2$ に比例するので基準運動量の粒子ではゼロになり，薄肉レンズ極限でも
ゼロになる。$A$，$B$ は $\mathcal{O}((Lg)^4)$ である一方，写像自体は $L$ の次数だからである。}

\en{\textbf{Why the defect appears.} The exact map is symplectic, so the third-order terms $\mathbf{t}$ of the map
compensate the second-order defect,}%
\ja{\textbf{欠陥が現れる理由。}厳密な写像はシンプレクティックであるから，写像の 3 次の項 $\mathbf{t}$ が
2 次の欠陥を補償する。}
\begin{equation}
  \mathbf{W}^\mathsf{T}\mathbf{J}\mathbf{W}
  = -\big(\mathbf{M}^\mathsf{T}\mathbf{J}\mathbf{T} + \mathbf{T}^\mathsf{T}\mathbf{J}\mathbf{M}\big)
  + \mathcal{O}(u_0^3) ,
  \qquad \mathbf{T} = \frac{\partial\mathbf{t}}{\partial\mathbf{u}_0} .
  \label{eq:defect_compensation}
\end{equation}
\en{The defect is therefore not an error of the coefficients $T_{ijk}$ but a consequence of truncating at second
order.}%
\ja{したがって欠陥は係数 $T_{ijk}$ の誤りではなく，2 次で打ち切ったことの結果である。}

\en{\textbf{Generating function.} A symplectic map has a generating function. For a type-2 function
$G(\mathbf{q}_0, \mathbf{p}_0)$ the second-order part must satisfy $\mathbf{w}_q = \partial G/\partial\mathbf{p}_0$ and
$\mathbf{w}_p = -\partial G/\partial\mathbf{q}_0$, which requires}%
\ja{\textbf{生成関数。}シンプレクティック写像は生成関数を持つ。2 型の関数 $G(\mathbf{q}_0, \mathbf{p}_0)$ に対して
2 次の部分は $\mathbf{w}_q = \partial G/\partial\mathbf{p}_0$，$\mathbf{w}_p = -\partial G/\partial\mathbf{q}_0$ を
満たす必要があり，そのためには}
\begin{equation}
  \frac{\partial w_{q,i}}{\partial q_{0j}} + \frac{\partial w_{p,j}}{\partial p_{0i}} = 0
  \label{eq:integrability}
\end{equation}
\en{for all $i, j$. The drift space satisfies \eqref{eq:integrability} with $G = L\,K$ of \eqref{eq:K_drift}. The
sector dipole magnet violates it in one entry and the quadrupole magnet in four, which is an independent proof that
their second-order maps are not symplectic.}%
\ja{がすべての $i$，$j$ について成り立つ必要がある。ドリフトスペースは \eqref{eq:K_drift} の $G = L\,K$ により
\eqref{eq:integrability} を満たす。扇型偏向磁石は 1 箇所，四極磁石は 4 箇所でこれを破り，これらの 2 次の写像が
シンプレクティックでないことの独立な証明になる。}

\en{\textbf{Practical consequence.} Composing second-order maps of many elements accumulates the defect, because
each map contributes an $\mathcal{O}(u_0^2)$ error that is not cancelled by the next one. For long-term tracking or
for a ring, either keep the third-order terms or use a symplectic integrator, which splits the Hamiltonian into
exactly solvable symplectic pieces instead of truncating the map.}%
\ja{\textbf{実用上の帰結。}多数の要素の 2 次の写像を合成すると欠陥が蓄積する。各写像が
$\mathcal{O}(u_0^2)$ の誤差を持ち，次の写像では相殺されないためである。長期のトラッキングやリングに対しては，
3 次の項を残すか，写像を打ち切る代わりにハミルトニアンを厳密に解けるシンプレクティックな部分に分割する
シンプレクティック積分器を用いる必要がある。}

% Slides for this section.
\slidetitle{\en{Symplecticity of the truncated map}\ja{打ち切った写像のシンプレクティック性}}
\slidemath{\mathbf{D}^\mathsf{T}\mathbf{J}\mathbf{D} - \mathbf{J} = \big(\mathbf{M}^\mathsf{T}\mathbf{J}\mathbf{W} + \mathbf{W}^\mathsf{T}\mathbf{J}\mathbf{M}\big) + \mathbf{W}^\mathsf{T}\mathbf{J}\mathbf{W}}
\slidepoint{\en{The O(u_0) group vanishes; the first defect is W^T J W, quadratic in u_0}%
            \ja{O(u_0) の群はゼロ。最初の欠陥は W^T J W で u_0 の 2 次}}
\slidepoint{\en{Drift: exact (flow of a polynomial Hamiltonian); dipole: sin^2(theta); quadrupole: delta_0^2}%
            \ja{ドリフト: 厳密（多項式ハミルトニアンの流れ），偏向: sin^2(theta)，四極: delta_0^2}}
\slidepoint{\en{The third-order terms compensate the defect: truncation, not a coefficient error}%
            \ja{3 次の項が欠陥を補償する。係数の誤りではなく打ち切りの結果}}

\section{\en{Checks}\ja{確認}}

\en{The notebook performs the following checks for every element, including the quadrupole with $k < 0$.}%
\ja{ノートブックでは，$k < 0$ の四極磁石を含む各要素について次の確認を行っている。}
\begin{enumerate}
  \item \en{Duhamel's integrals agree with the closed forms \eqref{eq:w_drift}--\eqref{eq:w_quad}.}%
        \ja{Duhamel の積分は閉じた式 \eqref{eq:w_drift}--\eqref{eq:w_quad} と一致する。}
  \item \en{The closed forms satisfy \eqref{eq:w_ode} and $\mathbf{w}(0) = 0$ (residual of the order of $10^{-16}$).}%
        \ja{閉じた式は \eqref{eq:w_ode} と $\mathbf{w}(0) = 0$ を満たす（残差は $10^{-16}$ 程度）。}
  \item \en{Second-order symplectic condition. The map $\mathbf{P}(\mathbf{u}_0) = \mathbf{M}\mathbf{u}_0 +
    \mathbf{w}(\mathbf{u}_0)$ is symplectic through first order in $\mathbf{u}_0$ if
    $\mathbf{M}^\mathsf{T}\mathbf{J}\mathbf{W} + \mathbf{W}^\mathsf{T}\mathbf{J}\mathbf{M} = 0$ with
    $\mathbf{W} = \partial\mathbf{w}/\partial\mathbf{u}_0$. This relates elements of $T_{ijk}$ and $R_{ij}$, and it
    holds for all elements.}%
        \ja{2 次のシンプレクティック条件。写像 $\mathbf{P}(\mathbf{u}_0) = \mathbf{M}\mathbf{u}_0 +
        \mathbf{w}(\mathbf{u}_0)$ が $\mathbf{u}_0$ の 1 次までシンプレクティックであるための条件は，
        $\mathbf{W} = \partial\mathbf{w}/\partial\mathbf{u}_0$ として
        $\mathbf{M}^\mathsf{T}\mathbf{J}\mathbf{W} + \mathbf{W}^\mathsf{T}\mathbf{J}\mathbf{M} = 0$ である。
        これは $T_{ijk}$ と $R_{ij}$ の要素を結びつけ，すべての要素で成り立つ。}
  \item \en{Composition: for two sections of the same element,
    $\mathbf{w}(L_1 + L_2) = \mathbf{M}(L_2)\,\mathbf{w}(L_1) + \mathbf{w}_{L_2}(\mathbf{M}(L_1)\,\mathbf{u}_0)$.}%
        \ja{合成則: 同じ要素の二つの区間について，
        $\mathbf{w}(L_1 + L_2) = \mathbf{M}(L_2)\,\mathbf{w}(L_1) + \mathbf{w}_{L_2}(\mathbf{M}(L_1)\,\mathbf{u}_0)$
        が成り立つ。}
  \item \en{Limits: the dipole tends to the drift space for $h \to 0$ and the quadrupole for $g \to 0$.}%
        \ja{極限: $h \to 0$ で偏向磁石が，$g \to 0$ で四極磁石がドリフトスペースに近づく。}
  \item \en{Comparison with the exact EOM, integrated numerically for the amplitude $a = 2\times 10^{-3}$ and for
    $a/2$: the error of the linear map scales as $a^2$ (ratio 4.0), and that of the second-order map as $a^3$
    (ratio 8.0), for every element. At $a = 2\times 10^{-3}$ the second-order map is more accurate than the
    linear map by a factor of 300 to 400.}%
        \ja{厳密な運動方程式との比較。振幅 $a = 2\times 10^{-3}$ と $a/2$ で数値積分すると，すべての要素で，
        線形写像の誤差は $a^2$（比 4.0），2 次の写像の誤差は $a^3$（比 8.0）に比例する。$a = 2\times 10^{-3}$ では，
        2 次の写像は線形写像より 300 倍から 400 倍精度が高い。}
\end{enumerate}

% Slides for this section.
\slidetitle{\en{Checks}\ja{確認}}
\slidepoint{\en{Duhamel integrals match the closed forms; ODE residual about 1e-16}%
            \ja{Duhamel の積分は閉じた式と一致，ODE の残差は 1e-16 程度}}
\slidepoint{\en{Second-order symplectic condition and composition hold for every element}%
            \ja{2 次のシンプレクティック条件と合成則が全要素で成立}}
\slidepoint{\en{Error scaling against the exact EOM: a^2 (linear), a^3 (second order)}%
            \ja{厳密な運動方程式との誤差: 線形は a^2，2 次は a^3}}

\section{\en{Summary and open points}\ja{まとめと未解決の点}}

\begin{enumerate}
  \item \en{$\mathbf{w}(L) = \int_0^L \mathbf{M}(L-s)\,\mathbf{f}_2(\mathbf{M}(s)\mathbf{u}_0)\,ds$, and $T_{ijk}$ is the
    coefficient of $u_{0j}u_{0k}$ ($j \le k$) in $w_i$.}%
        \ja{$\mathbf{w}(L) = \int_0^L \mathbf{M}(L-s)\,\mathbf{f}_2(\mathbf{M}(s)\mathbf{u}_0)\,ds$ であり，
        $T_{ijk}$ は $w_i$ における $u_{0j}u_{0k}$（$j \le k$）の係数である。}
  \item \en{Drift space: \eqref{eq:w_drift}, 5 nonzero elements. Sector dipole magnet: \eqref{eq:w_bend}, 18 nonzero
    elements. Quadrupole magnet: \eqref{eq:w_quad}, 15 nonzero elements.}%
        \ja{ドリフトスペース: \eqref{eq:w_drift}，非零の要素は 5 個。扇型偏向磁石: \eqref{eq:w_bend}，18 個。
        四極磁石: \eqref{eq:w_quad}，15 個。}
  \item \en{A thin quadrupole has no second-order terms in the canonical variables.}%
        \ja{薄い四極磁石は，正準変数では 2 次の項を持たない。}
  \item \en{Symplecticity: the map truncated at second order is symplectic through first order in $\mathbf{u}_0$ only.
    The first defect is $\mathbf{W}^\mathsf{T}\mathbf{J}\mathbf{W}$, quadratic in $\mathbf{u}_0$. It vanishes
    identically for the drift space, scales as $\sin^2\theta$ for the sector dipole magnet, and as $\delta_0^2$ for
    the quadrupole magnet. The third-order terms of the exact map compensate it.}%
        \ja{シンプレクティック性: 2 次で打ち切った写像は $\mathbf{u}_0$ の 1 次までしかシンプレクティックでない。
        最初の欠陥は $\mathbf{W}^\mathsf{T}\mathbf{J}\mathbf{W}$ で $\mathbf{u}_0$ の 2 次である。ドリフトスペースでは
        恒等的にゼロ，扇型偏向磁石では $\sin^2\theta$，四極磁石では $\delta_0^2$ に比例する。厳密な写像の 3 次の項が
        これを補償する。}
  \item \en{Open points, not covered here: the chromaticity of a ring from these maps (it needs a periodic ring and
    the dispersion) and the conversion of $T_{ijk}$ to the slopes $x'$ and $y'$; the third-order map and its
    comparison with a symplectic integrator; fringe fields, edge focusing, and the combined-function magnet;
    sextupole and octupole terms, quadrupole curvature corrections, the skew quadrupole, and RF.}%
        \ja{本メモで扱わない点: これらの写像から得られるリングのクロマティシティ（周期的なリングと分散が必要）と，
        $T_{ijk}$ から傾き $x'$，$y'$ への変換，3 次の写像とシンプレクティック積分器との比較，端部の漏れ磁場，
        エッジ集束，複合型磁石，六極磁石と八極磁石の項，四極磁石の曲率補正，スキュー四極磁石，RF。}
\end{enumerate}

% Slides for this section.
\slidetitle{\en{Summary}\ja{まとめ}}
\slidemath{\mathbf{w}(L) = \int_0^L \mathbf{M}(L-s)\,\mathbf{f}_2\big(\mathbf{M}(s)\mathbf{u}_0\big)\,ds}
\slidepoint{\en{Nonzero T_ijk: drift 5, sector dipole 18, quadrupole 15}%
            \ja{非零の T_ijk: ドリフト 5 個，偏向磁石 18 個，四極磁石 15 個}}
\slidepoint{\en{Second-order map: symplectic only through first order; defect W^T J W}%
            \ja{2 次の写像は 1 次までしかシンプレクティックでない。欠陥は W^T J W}}
\slidepoint{\en{Not covered: chromaticity of a ring, edge focusing, sextupoles, RF}%
            \ja{未対応: リングのクロマティシティ，エッジ集束，六極磁石，RF}}

\printbibliography[title={\en{References}\ja{参考文献}}]

\section*{\en{License}\ja{ライセンス}}
\en{This document is licensed under CC BY 4.0. Code snippets are GPL-3.0-or-later. See \texttt{LICENSING.md}.}%
\ja{本文書は CC BY 4.0 で提供する。コードの断片は GPL-3.0-or-later である。詳細は \texttt{LICENSING.md} を参照すること。}

\end{document}
